
\documentclass{article}
\usepackage{colortbl}
\usepackage{makecell}
\usepackage{multirow}
\usepackage{supertabular}

\begin{document}

\newcounter{utterance}

\twocolumn

{ \footnotesize  \setcounter{utterance}{1}
\setlength{\tabcolsep}{0pt}
\begin{supertabular}{c@{$\;$}|p{.15\linewidth}@{}p{.15\linewidth}p{.15\linewidth}p{.15\linewidth}p{.15\linewidth}p{.15\linewidth}}

    \# & $\;$A & \multicolumn{4}{c}{Game Master} & $\;\:$B\\
    \hline 

    \theutterance \stepcounter{utterance}  

    & & \multicolumn{4}{p{0.6\linewidth}}{\cellcolor[rgb]{0.9,0.9,0.9}{%
	\makecell[{{p{\linewidth}}}]{% 
	  \tt {\tiny [A$\langle$GM]}  
	 State giocando a un gioco di negoziazione in cui dovete accordarvi su come dividere una serie di oggetti tra voi e un altro giocatore.\\ \tt \\ \tt Regole:\\ \tt (a) A Lei e all'altro giocatore viene dato un insieme di oggetti. Ognuno di voi riceve anche una funzione di valore segreta, che rappresenta il valore di ciascun tipo di oggetto.\\ \tt (b) Si scambiano messaggi con l'altro giocatore per concordare chi si aggiudica gli oggetti. Potete inviare un massimo di 5 messaggi ciascuno, oppure terminare in anticipo facendo una proposta segreta in qualsiasi momento.\\ \tt (c) A ciascuno di voi viene chiesto di inviare una proposta segreta indicando gli oggetti che desiderate, formattata tra parentesi quadre come segue: "[Proposta: <numero> <nome oggetto>, <numero> <nome oggetto>, <...>]".\\ \tt (d) Se le vostre proposte sono complementari, cioè ci sono abbastanza oggetti per soddisfare entrambe le proposte, a ciascun giocatore viene assegnato un punteggio basato sulla somma dei valori degli oggetti ricevuti. In caso contrario, entrambi ricevono zero punti.\\ \tt (e) L'obiettivo è massimizzare il proprio punteggio.\\ \tt \\ \tt Cominciamo.\\ \tt \\ \tt L'insieme degli oggetti disponibili è:\\ \tt \\ \tt 3 pastelli, 1 gesso, 2 bulloni\\ \tt \\ \tt La funzione valore segreta è:\\ \tt \\ \tt pastello: 0, gesso: 0, bullone: 5\\ \tt \\ \tt IMPORTANTE: i vostri messaggi, a meno che non si tratti di una proposta segreta, vengono trasmessi direttamente all'altro giocatore, quindi non includete alcuna risposta alle regole o testo di annuncio del vostro messaggio. Per fare una proposta segreta, utilizzate il formato indicato. Non utilizzare delle parentesi quadre quando si comunica all'altro giocatore, altrimenti verrà interpretata come una proposta segreta.\\ \tt  
	  } 
	   } 
	   } 
	 & \\ 
 

    \theutterance \stepcounter{utterance}  

    & \multicolumn{4}{p{0.6\linewidth}}{\cellcolor[rgb]{0.8,1,0.9}{%
	\makecell[{{p{\linewidth}}}]{% 
	  \tt {\tiny [A$\rangle$GM]}  
	 Ciao, vorrei discutere di come dividere gli oggetti. Hai una preferenza per alcuni di essi? 
	  } 
	   } 
	   } 
	 & & \\ 
 

    \theutterance \stepcounter{utterance}  

    & & & \multicolumn{2}{p{0.3\linewidth}}{\cellcolor[rgb]{0.95,0.95,0.95}{%
	\makecell[{{p{\linewidth}}}]{% 
	  \tt {\tiny [GM$|$GM]}  
	 continue 
	  } 
	   } 
	   } 
	 & & \\ 
 

    \theutterance \stepcounter{utterance}  

    & & & \multicolumn{2}{p{0.3\linewidth}}{\cellcolor[rgb]{0.95,0.95,0.95}{%
	\makecell[{{p{\linewidth}}}]{% 
	  \tt {\tiny [GM$|$GM]}  
	 Ciao, vorrei discutere di come dividere gli oggetti. Hai una preferenza per alcuni di essi? 
	  } 
	   } 
	   } 
	 & & \\ 
 

    \theutterance \stepcounter{utterance}  

    & & \multicolumn{4}{p{0.6\linewidth}}{\cellcolor[rgb]{0.9,0.9,0.9}{%
	\makecell[{{p{\linewidth}}}]{% 
	  \tt {\tiny [GM$\rangle$B]}  
	 State giocando a un gioco di negoziazione in cui dovete accordarvi su come dividere una serie di oggetti tra voi e un altro giocatore.\\ \tt \\ \tt Regole:\\ \tt (a) A Lei e all'altro giocatore viene dato un insieme di oggetti. Ognuno di voi riceve anche una funzione di valore segreta, che rappresenta il valore di ciascun tipo di oggetto.\\ \tt (b) Si scambiano messaggi con l'altro giocatore per concordare chi si aggiudica gli oggetti. Potete inviare un massimo di 5 messaggi ciascuno, oppure terminare in anticipo facendo una proposta segreta in qualsiasi momento.\\ \tt (c) A ciascuno di voi viene chiesto di inviare una proposta segreta indicando gli oggetti che desiderate, formattata tra parentesi quadre come segue: "[Proposta: <numero> <nome oggetto>, <numero> <nome oggetto>, <...>]".\\ \tt (d) Se le vostre proposte sono complementari, cioè ci sono abbastanza oggetti per soddisfare entrambe le proposte, a ciascun giocatore viene assegnato un punteggio basato sulla somma dei valori degli oggetti ricevuti. In caso contrario, entrambi ricevono zero punti.\\ \tt (e) L'obiettivo è massimizzare il proprio punteggio.\\ \tt \\ \tt Cominciamo.\\ \tt \\ \tt L'insieme degli oggetti disponibili è:\\ \tt \\ \tt 3 pastelli, 1 gesso, 2 bulloni\\ \tt \\ \tt La funzione valore segreta è:\\ \tt \\ \tt pastello: 1, gesso: 3, bullone: 2\\ \tt \\ \tt IMPORTANTE: i vostri messaggi, a meno che non si tratti di una proposta segreta, vengono trasmessi direttamente all'altro giocatore, quindi non includete alcuna risposta alle regole o testo di annuncio del vostro messaggio. Per fare una proposta segreta, utilizzate il formato indicato. Non utilizzare delle parentesi quadre quando si comunica all'altro giocatore, altrimenti verrà interpretata come una proposta segreta.\\ \tt \\ \tt \\ \tt Ciao, vorrei discutere di come dividere gli oggetti. Hai una preferenza per alcuni di essi? 
	  } 
	   } 
	   } 
	 & \\ 
 

    \theutterance \stepcounter{utterance}  

    & & & \multicolumn{4}{p{0.6\linewidth}}{\cellcolor[rgb]{1,0.85,0.72}{%
	\makecell[{{p{\linewidth}}}]{% 
	  \tt {\tiny [GM$\langle$B]}  
	 Ciao! Sì, ho una certa preferenza per i bulloni, ma sono anche interessato ai pastelli. Il gesso mi sembra meno importante, ma potrebbe essere comunque utile. E tu, quali oggetti ti interessano di più? 
	  } 
	   } 
	   } 
	  \\ 
 

    \theutterance \stepcounter{utterance}  

    & & & \multicolumn{2}{p{0.3\linewidth}}{\cellcolor[rgb]{0.95,0.95,0.95}{%
	\makecell[{{p{\linewidth}}}]{% 
	  \tt {\tiny [GM$|$GM]}  
	 continue 
	  } 
	   } 
	   } 
	 & & \\ 
 

    \theutterance \stepcounter{utterance}  

    & & & \multicolumn{2}{p{0.3\linewidth}}{\cellcolor[rgb]{0.95,0.95,0.95}{%
	\makecell[{{p{\linewidth}}}]{% 
	  \tt {\tiny [GM$|$GM]}  
	 Ciao! Sì, ho una certa preferenza per i bulloni, ma sono anche interessato ai pastelli. Il gesso mi sembra meno importante, ma potrebbe essere comunque utile. E tu, quali oggetti ti interessano di più? 
	  } 
	   } 
	   } 
	 & & \\ 
 

    \theutterance \stepcounter{utterance}  

    & & \multicolumn{4}{p{0.6\linewidth}}{\cellcolor[rgb]{0.9,0.9,0.9}{%
	\makecell[{{p{\linewidth}}}]{% 
	  \tt {\tiny [A$\langle$GM]}  
	 Ciao! Sì, ho una certa preferenza per i bulloni, ma sono anche interessato ai pastelli. Il gesso mi sembra meno importante, ma potrebbe essere comunque utile. E tu, quali oggetti ti interessano di più? 
	  } 
	   } 
	   } 
	 & \\ 
 

    \theutterance \stepcounter{utterance}  

    & \multicolumn{4}{p{0.6\linewidth}}{\cellcolor[rgb]{0.8,1,0.9}{%
	\makecell[{{p{\linewidth}}}]{% 
	  \tt {\tiny [A$\rangle$GM]}  
	 Anch'io ritengo che i bulloni siano molto importanti, ma sono aperto a trovare un accordo che ci soddisfi entrambi. I pastelli potrebbero essere una buona opzione per te, mentre io sono più interessato ai bulloni. Il gesso potrebbe essere un oggetto di scambio o una concessione. Che ne pensi di una divisione in cui prendiamo entrambi un bullone e poi decidiamo cosa fare con il resto? 
	  } 
	   } 
	   } 
	 & & \\ 
 

    \theutterance \stepcounter{utterance}  

    & & & \multicolumn{2}{p{0.3\linewidth}}{\cellcolor[rgb]{0.95,0.95,0.95}{%
	\makecell[{{p{\linewidth}}}]{% 
	  \tt {\tiny [GM$|$GM]}  
	 continue 
	  } 
	   } 
	   } 
	 & & \\ 
 

    \theutterance \stepcounter{utterance}  

    & & & \multicolumn{2}{p{0.3\linewidth}}{\cellcolor[rgb]{0.95,0.95,0.95}{%
	\makecell[{{p{\linewidth}}}]{% 
	  \tt {\tiny [GM$|$GM]}  
	 Anch'io ritengo che i bulloni siano molto importanti, ma sono aperto a trovare un accordo che ci soddisfi entrambi. I pastelli potrebbero essere una buona opzione per te, mentre io sono più interessato ai bulloni. Il gesso potrebbe essere un oggetto di scambio o una concessione. Che ne pensi di una divisione in cui prendiamo entrambi un bullone e poi decidiamo cosa fare con il resto? 
	  } 
	   } 
	   } 
	 & & \\ 
 

    \theutterance \stepcounter{utterance}  

    & & \multicolumn{4}{p{0.6\linewidth}}{\cellcolor[rgb]{0.9,0.9,0.9}{%
	\makecell[{{p{\linewidth}}}]{% 
	  \tt {\tiny [GM$\rangle$B]}  
	 Anch'io ritengo che i bulloni siano molto importanti, ma sono aperto a trovare un accordo che ci soddisfi entrambi. I pastelli potrebbero essere una buona opzione per te, mentre io sono più interessato ai bulloni. Il gesso potrebbe essere un oggetto di scambio o una concessione. Che ne pensi di una divisione in cui prendiamo entrambi un bullone e poi decidiamo cosa fare con il resto? 
	  } 
	   } 
	   } 
	 & \\ 
 

    \theutterance \stepcounter{utterance}  

    & & & \multicolumn{4}{p{0.6\linewidth}}{\cellcolor[rgb]{1,0.85,0.72}{%
	\makecell[{{p{\linewidth}}}]{% 
	  \tt {\tiny [GM$\langle$B]}  
	 Mi sembra una buona base di partenza. Prendere un bullone ciascuno semplificherebbe le cose. Per quanto riguarda il resto, potremmo dividerci i pastelli, visto che ne hai una certa disponibilità. Il gesso potrebbe essere assegnato a una delle due parti in base a ulteriori negoziati. Come pensi di dividere esattamente i pastelli? Una divisione equa potrebbe essere due pastelli per te e uno per me, visto che ce ne sono tre in totale? 
	  } 
	   } 
	   } 
	  \\ 
 

    \theutterance \stepcounter{utterance}  

    & & & \multicolumn{2}{p{0.3\linewidth}}{\cellcolor[rgb]{0.95,0.95,0.95}{%
	\makecell[{{p{\linewidth}}}]{% 
	  \tt {\tiny [GM$|$GM]}  
	 continue 
	  } 
	   } 
	   } 
	 & & \\ 
 

    \theutterance \stepcounter{utterance}  

    & & & \multicolumn{2}{p{0.3\linewidth}}{\cellcolor[rgb]{0.95,0.95,0.95}{%
	\makecell[{{p{\linewidth}}}]{% 
	  \tt {\tiny [GM$|$GM]}  
	 Mi sembra una buona base di partenza. Prendere un bullone ciascuno semplificherebbe le cose. Per quanto riguarda il resto, potremmo dividerci i pastelli, visto che ne hai una certa disponibilità. Il gesso potrebbe essere assegnato a una delle due parti in base a ulteriori negoziati. Come pensi di dividere esattamente i pastelli? Una divisione equa potrebbe essere due pastelli per te e uno per me, visto che ce ne sono tre in totale? 
	  } 
	   } 
	   } 
	 & & \\ 
 

    \theutterance \stepcounter{utterance}  

    & & \multicolumn{4}{p{0.6\linewidth}}{\cellcolor[rgb]{0.9,0.9,0.9}{%
	\makecell[{{p{\linewidth}}}]{% 
	  \tt {\tiny [A$\langle$GM]}  
	 Mi sembra una buona base di partenza. Prendere un bullone ciascuno semplificherebbe le cose. Per quanto riguarda il resto, potremmo dividerci i pastelli, visto che ne hai una certa disponibilità. Il gesso potrebbe essere assegnato a una delle due parti in base a ulteriori negoziati. Come pensi di dividere esattamente i pastelli? Una divisione equa potrebbe essere due pastelli per te e uno per me, visto che ce ne sono tre in totale? 
	  } 
	   } 
	   } 
	 & \\ 
 

    \theutterance \stepcounter{utterance}  

    & \multicolumn{4}{p{0.6\linewidth}}{\cellcolor[rgb]{0.8,1,0.9}{%
	\makecell[{{p{\linewidth}}}]{% 
	  \tt {\tiny [A$\rangle$GM]}  
	 La tua proposta sui pastelli sembra ragionevole, considerando che ci sono tre pastelli in totale. Due per me e uno per te potrebbe essere una soluzione accettabile. Tuttavia, vorrei chiarire la questione del gesso. Dal momento che a entrambi sono stati assegnati un bullone, potremmo considerare di assegnare il gesso a te, in cambio del fatto che io prendo due pastelli e tu ne prendi uno. Questo modo di procedere ti sembra equo? 
	  } 
	   } 
	   } 
	 & & \\ 
 

    \theutterance \stepcounter{utterance}  

    & & & \multicolumn{2}{p{0.3\linewidth}}{\cellcolor[rgb]{0.95,0.95,0.95}{%
	\makecell[{{p{\linewidth}}}]{% 
	  \tt {\tiny [GM$|$GM]}  
	 continue 
	  } 
	   } 
	   } 
	 & & \\ 
 

    \theutterance \stepcounter{utterance}  

    & & & \multicolumn{2}{p{0.3\linewidth}}{\cellcolor[rgb]{0.95,0.95,0.95}{%
	\makecell[{{p{\linewidth}}}]{% 
	  \tt {\tiny [GM$|$GM]}  
	 La tua proposta sui pastelli sembra ragionevole, considerando che ci sono tre pastelli in totale. Due per me e uno per te potrebbe essere una soluzione accettabile. Tuttavia, vorrei chiarire la questione del gesso. Dal momento che a entrambi sono stati assegnati un bullone, potremmo considerare di assegnare il gesso a te, in cambio del fatto che io prendo due pastelli e tu ne prendi uno. Questo modo di procedere ti sembra equo? 
	  } 
	   } 
	   } 
	 & & \\ 
 

    \theutterance \stepcounter{utterance}  

    & & \multicolumn{4}{p{0.6\linewidth}}{\cellcolor[rgb]{0.9,0.9,0.9}{%
	\makecell[{{p{\linewidth}}}]{% 
	  \tt {\tiny [GM$\rangle$B]}  
	 La tua proposta sui pastelli sembra ragionevole, considerando che ci sono tre pastelli in totale. Due per me e uno per te potrebbe essere una soluzione accettabile. Tuttavia, vorrei chiarire la questione del gesso. Dal momento che a entrambi sono stati assegnati un bullone, potremmo considerare di assegnare il gesso a te, in cambio del fatto che io prendo due pastelli e tu ne prendi uno. Questo modo di procedere ti sembra equo? 
	  } 
	   } 
	   } 
	 & \\ 
 

    \theutterance \stepcounter{utterance}  

    & & & \multicolumn{4}{p{0.6\linewidth}}{\cellcolor[rgb]{1,0.85,0.72}{%
	\makecell[{{p{\linewidth}}}]{% 
	  \tt {\tiny [GM$\langle$B]}  
	 Mi sembra una proposta molto equa. Assegnare il gesso a me in cambio del fatto che tu prendi due pastelli e io ne prendo uno, dopo che abbiamo già concordato di prendere un bullone ciascuno, sembra un buon compromesso. Sembra che stiamo arrivando a un accordo che ci soddisfi entrambi. Prima di procedere con una proposta segreta formale, vorresti che ribadissi l'accordo in modo chiaro, così siamo certi di essere sulla stessa lunghezza d'onda? 
	  } 
	   } 
	   } 
	  \\ 
 

    \theutterance \stepcounter{utterance}  

    & & & \multicolumn{2}{p{0.3\linewidth}}{\cellcolor[rgb]{0.95,0.95,0.95}{%
	\makecell[{{p{\linewidth}}}]{% 
	  \tt {\tiny [GM$|$GM]}  
	 continue 
	  } 
	   } 
	   } 
	 & & \\ 
 

    \theutterance \stepcounter{utterance}  

    & & & \multicolumn{2}{p{0.3\linewidth}}{\cellcolor[rgb]{0.95,0.95,0.95}{%
	\makecell[{{p{\linewidth}}}]{% 
	  \tt {\tiny [GM$|$GM]}  
	 Mi sembra una proposta molto equa. Assegnare il gesso a me in cambio del fatto che tu prendi due pastelli e io ne prendo uno, dopo che abbiamo già concordato di prendere un bullone ciascuno, sembra un buon compromesso. Sembra che stiamo arrivando a un accordo che ci soddisfi entrambi. Prima di procedere con una proposta segreta formale, vorresti che ribadissi l'accordo in modo chiaro, così siamo certi di essere sulla stessa lunghezza d'onda? 
	  } 
	   } 
	   } 
	 & & \\ 
 

    \theutterance \stepcounter{utterance}  

    & & \multicolumn{4}{p{0.6\linewidth}}{\cellcolor[rgb]{0.9,0.9,0.9}{%
	\makecell[{{p{\linewidth}}}]{% 
	  \tt {\tiny [A$\langle$GM]}  
	 Mi sembra una proposta molto equa. Assegnare il gesso a me in cambio del fatto che tu prendi due pastelli e io ne prendo uno, dopo che abbiamo già concordato di prendere un bullone ciascuno, sembra un buon compromesso. Sembra che stiamo arrivando a un accordo che ci soddisfi entrambi. Prima di procedere con una proposta segreta formale, vorresti che ribadissi l'accordo in modo chiaro, così siamo certi di essere sulla stessa lunghezza d'onda? 
	  } 
	   } 
	   } 
	 & \\ 
 

    \theutterance \stepcounter{utterance}  

    & \multicolumn{4}{p{0.6\linewidth}}{\cellcolor[rgb]{0.8,1,0.9}{%
	\makecell[{{p{\linewidth}}}]{% 
	  \tt {\tiny [A$\rangle$GM]}  
	 Sì, sarebbe utile ribadire l'accordo in modo chiaro per assicurarci di essere d'accordo su tutti i punti. Secondo me, l'accordo attuale prevede che entrambi prendiamo un bullone, poi tu prendi uno dei pastelli e io prendo i restanti due pastelli, e infine il gesso va a te. Se questo è corretto, potremmo procedere a inviare le nostre proposte segrete formalizzando questo accordo. Se falta qualcosa o c'è stato un malinteso, è meglio chiarirlo adesso. 
	  } 
	   } 
	   } 
	 & & \\ 
 

    \theutterance \stepcounter{utterance}  

    & & & \multicolumn{2}{p{0.3\linewidth}}{\cellcolor[rgb]{0.95,0.95,0.95}{%
	\makecell[{{p{\linewidth}}}]{% 
	  \tt {\tiny [GM$|$GM]}  
	 continue 
	  } 
	   } 
	   } 
	 & & \\ 
 

    \theutterance \stepcounter{utterance}  

    & & & \multicolumn{2}{p{0.3\linewidth}}{\cellcolor[rgb]{0.95,0.95,0.95}{%
	\makecell[{{p{\linewidth}}}]{% 
	  \tt {\tiny [GM$|$GM]}  
	 Sì, sarebbe utile ribadire l'accordo in modo chiaro per assicurarci di essere d'accordo su tutti i punti. Secondo me, l'accordo attuale prevede che entrambi prendiamo un bullone, poi tu prendi uno dei pastelli e io prendo i restanti due pastelli, e infine il gesso va a te. Se questo è corretto, potremmo procedere a inviare le nostre proposte segrete formalizzando questo accordo. Se falta qualcosa o c'è stato un malinteso, è meglio chiarirlo adesso. 
	  } 
	   } 
	   } 
	 & & \\ 
 

    \theutterance \stepcounter{utterance}  

    & & \multicolumn{4}{p{0.6\linewidth}}{\cellcolor[rgb]{0.9,0.9,0.9}{%
	\makecell[{{p{\linewidth}}}]{% 
	  \tt {\tiny [GM$\rangle$B]}  
	 Sì, sarebbe utile ribadire l'accordo in modo chiaro per assicurarci di essere d'accordo su tutti i punti. Secondo me, l'accordo attuale prevede che entrambi prendiamo un bullone, poi tu prendi uno dei pastelli e io prendo i restanti due pastelli, e infine il gesso va a te. Se questo è corretto, potremmo procedere a inviare le nostre proposte segrete formalizzando questo accordo. Se falta qualcosa o c'è stato un malinteso, è meglio chiarirlo adesso. 
	  } 
	   } 
	   } 
	 & \\ 
 

    \theutterance \stepcounter{utterance}  

    & & & \multicolumn{4}{p{0.6\linewidth}}{\cellcolor[rgb]{1,0.85,0.72}{%
	\makecell[{{p{\linewidth}}}]{% 
	  \tt {\tiny [GM$\langle$B]}  
	 Sì, l'accordo che hai descritto è esattamente quello a cui eravamo arrivati: un bullone a testa, due pastelli a te e uno a me, e il gesso a me. Sembra che abbiamo coperto tutti gli oggetti e che l'accordo sia equo per entrambi. Non vedo motivi per cui non procedere con le proposte segrete formalizzando questo accordo. Sono pronto a inviare la mia proposta segreta sulla base di quanto discusso. [Proposta: 1 bullone, 1 pastello, 1 gesso] 
	  } 
	   } 
	   } 
	  \\ 
 

    \theutterance \stepcounter{utterance}  

    & & & \multicolumn{2}{p{0.3\linewidth}}{\cellcolor[rgb]{0.95,0.95,0.95}{%
	\makecell[{{p{\linewidth}}}]{% 
	  \tt {\tiny [GM$|$GM]}  
	 proposal 
	  } 
	   } 
	   } 
	 & & \\ 
 

    \theutterance \stepcounter{utterance}  

    & & & \multicolumn{2}{p{0.3\linewidth}}{\cellcolor[rgb]{0.95,0.95,0.95}{%
	\makecell[{{p{\linewidth}}}]{% 
	  \tt {\tiny [GM$|$GM]}  
	 [1, 1, 1] 
	  } 
	   } 
	   } 
	 & & \\ 
 

    \theutterance \stepcounter{utterance}  

    & & \multicolumn{4}{p{0.6\linewidth}}{\cellcolor[rgb]{0.9,0.9,0.9}{%
	\makecell[{{p{\linewidth}}}]{% 
	  \tt {\tiny [A$\langle$GM]}  
	 L'altro giocatore ha presentato una proposta segreta. Ora tocca a Lei presentare una sua proposta.\\ \tt  
	  } 
	   } 
	   } 
	 & \\ 
 

    \theutterance \stepcounter{utterance}  

    & \multicolumn{4}{p{0.6\linewidth}}{\cellcolor[rgb]{0.8,1,0.9}{%
	\makecell[{{p{\linewidth}}}]{% 
	  \tt {\tiny [A$\rangle$GM]}  
	 [Proposta: 1 bullone, 2 pastelli] 
	  } 
	   } 
	   } 
	 & & \\ 
 

    \theutterance \stepcounter{utterance}  

    & & & \multicolumn{2}{p{0.3\linewidth}}{\cellcolor[rgb]{0.95,0.95,0.95}{%
	\makecell[{{p{\linewidth}}}]{% 
	  \tt {\tiny [GM$|$GM]}  
	 proposal 
	  } 
	   } 
	   } 
	 & & \\ 
 

    \theutterance \stepcounter{utterance}  

    & & & \multicolumn{2}{p{0.3\linewidth}}{\cellcolor[rgb]{0.95,0.95,0.95}{%
	\makecell[{{p{\linewidth}}}]{% 
	  \tt {\tiny [GM$|$GM]}  
	 [2, 0, 1] 
	  } 
	   } 
	   } 
	 & & \\ 
 

    \theutterance \stepcounter{utterance}  

    & & & \multicolumn{2}{p{0.3\linewidth}}{\cellcolor[rgb]{0.95,0.95,0.95}{%
	\makecell[{{p{\linewidth}}}]{% 
	  \tt {\tiny [GM$|$GM]}  
	 [[2, 0, 1], [1, 1, 1]] 
	  } 
	   } 
	   } 
	 & & \\ 
 

\end{supertabular}
}

\end{document}
